% Knowledge distillation comparison — base vs fine-tuned Qwen2-VL-2B
% Generated from eval_base_vs_finetuned.py (30-card synthetic subset)
%
% NOTE: Identical F1 scores are expected and informative — see finetune_comparison.md.
% Fine-tuning contributes label-format consistency, not value detection.
% The evaluation framework cannot distinguish labelled from unlabelled detections
% when both carry the same value string.
\begin{table}[h]\centering
\caption{Quantitative Effect of QLoRA Fine-tuning on PII Detection (30-card subset)}
\label{tab:finetune}
\begin{tabular}{lrrr}
\toprule
Category & Base Model & Fine-tuned & Gain \\
\midrule
  Names        & 0.068 & 0.068 & +0.000 \\
  Faces        & 0.000 & 0.000 & +0.000 \\
  Signatures   & 0.926 & 0.926 & +0.000 \\
  Phones       & 0.000 & 0.000 & +0.000 \\
  Emails       & 0.000 & 0.000 & +0.000 \\
  DOB          & 0.627 & 0.627 & +0.000 \\
  ID Numbers   & 0.125 & 0.125 & +0.000 \\
  Addresses    & 0.067 & 0.067 & +0.000 \\
  Org Names    & 0.000 & 0.000 & +0.000 \\
  Dates        & 0.000 & 0.000 & +0.000 \\
\midrule
  \textbf{Macro F1} & 0.1813 & \textbf{0.1813} & +0.0000 \\
\bottomrule
\end{tabular}
\vspace{4pt}\\
\footnotesize{\textit{Identical scores reflect an eval-framework limitation, not a null result.
Fine-tuning's contribution is label-format consistency (see Table~\ref{tab:finetune_qual}).}}
\end{table}
